%======================================================================
%  Using this companion -- unnumbered opening chapter
%======================================================================
\chapter*{Using this companion}\label{cc:intro}
\addcontentsline{toc}{chapter}{Using this companion}
\markboth{Using this companion}{Using this companion}

\framepart{Outline}

The lecture notes compute with Python: small scripts that print the numbers
quoted in the solutions. This companion does the same computations with the
finite element code Cast3M, chapter by chapter. Its chapters carry the numbers
of the chapters of the notes they accompany; Chapter~3, on the reciprocity
gap, has its own Cast3M files and is treated separately. Exercises are numbered
C\,\textit{chapter}.\textit{n}, so that Exercise~\ref{exo:cc-cylinder} is never confused with
Exercise~5.2 of the notes.

After this companion you should be able to
\begin{itemize}
\item check a result of the notes with the built-in procedures of Cast3M
  (\op{reso}, \op{pasapas});
\item program an integration algorithm of the notes directly on the fields of
  Cast3M: a return map on the Gauss-point fields, an equilibrium loop on the
  nodal fields;
\item drive Cast3M from outside, as the direct solver of an identification:
  from a shell script, from SciPy, from OpenTURNS and from JAX.
\end{itemize}

\section{Three ways to use Cast3M}\label{cc:sec-ways}
\index{Cast3M!three uses}%

\paragraph{The built-in procedures.}
\op{reso} solves a linear system with the constraints added by Lagrange
multipliers, and \op{pasapas} runs an incremental elastoplastic computation
with the return maps of the code. They check the closed forms of the notes:
the Kirsch factor, the reactions of a constrained spring chain, the stresses
in a stretched sheet, the residual stresses in a cylinder.

\paragraph{The algorithms of the notes, written in gibiane.}
Gibiane manipulates whole fields: a stress field at the Gauss points
(\op{MCHAML}), a displacement field at the nodes (\op{CHPOINT}). The radial
return of Box~5.2 (Section~\ref{B-sec:pn-j2}) is a sequence of operations on fields: a
trial stress, its von Mises norm, a mask of the plastic points, a scalar
Newton iteration done at all Gauss points at once, an update along the
deviator. The equilibrium loop of Box~5.3 is a residual and a linear solve.
Chapter~\ref{cc:ch5} writes both, and compares them with \op{pasapas}.

\paragraph{Cast3M as a black box.}
For an identification (Chapter~\ref{B-ch:id} of the notes) the direct solver
is called many times by a driver that only sees parameters in and numbers
out. A template file with placeholders, one directory per run, a CSV file for
the results: the same three lines serve a shell loop, SciPy, OpenTURNS and
JAX (Chapter~\ref{cc:ch7}).

\section{The exercises}\label{cc:sec-list}

\begin{gbox}{The Cast3M exercises and the notes}\label{cc:tab-list}
\small
\begin{tabular}{@{}l@{\hspace{3mm}}p{84mm}@{\hspace{3mm}}l@{}}
\textbf{ex.} & \textbf{what} & \textbf{notes}\\[1mm]
\ref{exo:cc-springs} & spring chain: $\mat{K}=\mat{B}\T\mat{C}\mat{B}$, multiplier, reciprocity
     & Sec.~\ref{B-sec:var-lagrange}, Ex.~\ref{B-exo:primal-dual}\\
\ref{exo:cc-kirsch} & Kirsch with four elements, balance of reactions & Ex.~\ref{B-exo:kirsch}\\
\ref{exo:cc-sheet} & sheet in plane strain, perfect plasticity & Ex.~\ref{B-exo:pl-sheet}\\
\ref{exo:cc-filament} & filament under a strain cycle & Ex.~\ref{B-exo:pl-cyclic}\\
\ref{exo:cc-residual} & cylinder: loading, unloading, residual stresses & Ex.~\ref{B-exo:pl-cyl}\\
\ref{exo:cc-rrj2} & radial return on Gauss-point fields & Ex.~\ref{B-exo:pn-shear}\\
\ref{exo:cc-cylinder} & cylinder: return map and equilibrium loop in gibiane & Ex.~\ref{B-exo:pn-fem}\\
\ref{exo:cc-bree} & Bree's tube, cycle by cycle & Ex.~\ref{B-exo:cy-bree}\\
\ref{exo:cc-id-q-scale}, \ref{exo:cc-id-q-step} & short questions: a load that scales, the step of finite differences & Ex.~\ref{B-exo:id-q-fd}\\
\ref{exo:cc-id-sweep} & the cost along one parameter, from the shell & Ex.~\ref{B-exo:id-truss}\\
\ref{exo:cc-id-scipy} & identification with SciPy & Ex.~\ref{B-exo:id-truss}\\
\ref{exo:cc-id-ot} & calibration with OpenTURNS & Ex.~\ref{B-exo:id-tools}\\
\ref{exo:cc-id-jax} & Cast3M inside a JAX program & Ex.~\ref{B-exo:id-tools}\\
\end{tabular}
\end{gbox}

\noindent The files live in the directory \texttt{cast3m/} of the notes, one
subdirectory per chapter. Each starts with a header that states the exercise,
the data and the reference values of the Python scripts; each prints its
results next to these values with \op{mess}. Run them with the local Cast3M
command, for instance \texttt{castem25 return\_map\_point.dgibi}.

\paragraph{Status.}
The gibiane files were written against the operator notices and the Cast3M
introduction, without a Cast3M installation, and a red note marks every one
that has not been run. The drivers of Chapter~\ref{cc:ch7} were run against a
stand-in that reads the same file and writes the same CSV with the Python model
of the notes; they reproduce the numbers of the book. A file goes into the notes
once it has been run with Cast3M.

\begin{gbox}{Constructions to check first}\label{cc:box-check}
\small
\begin{enumerate}[label=\textbf{\arabic*.},leftmargin=*,itemsep=1pt]
\item Arithmetic between two Gauss-point fields with the single component
  \texttt{SCAL} (\texttt{* / + -}), and \op{exco} renaming components both
  ways, \texttt{exco x 'SMXX' 'SCAL'} and \texttt{exco y 'SCAL' 'SMXX'}.
\item The component name \texttt{SCAL} of \op{vmis}; \op{masq} and \op{exp}
  on a Gauss-point field; the type of a stress field rebuilt from its
  components (\texttt{chan 'TYPE'} if an operator refuses it).
\item \texttt{chan 'CHAM' (coor 1 dom) mo 'STRESSES'}, the radius of the
  Gauss points; the name \texttt{EPSE} of the cumulated plastic strain.
\item Bar elements with \texttt{PLASTIQUE CINEMATIQUE}, and the meaning of
  its modulus \texttt{H} (Exercise~\ref{exo:cc-filament} decides it).
\item The \op{pasapas} entries \texttt{'PRECISION'} and
  \texttt{'TEMPERATURES'} with a \texttt{CHAR 'T'} loading; \op{sin} of a
  list of reals, in degrees.
\end{enumerate}
\end{gbox}

\framepart{Summary}

\begin{gbox*}
\begin{itemize}
\item \textbf{One chapter of the companion per chapter of the notes}, with
  the same number; exercises C\,\textit{chapter}.\textit{n}.
\item \textbf{Three uses of Cast3M}: its procedures check the notes; its
  fields carry the algorithms of the notes; as a black box it is the direct
  solver of an outside driver.
\item \textbf{Every file prints its own check}: the reference values of the
  Python scripts are in its header and next to its results.
\end{itemize}
\end{gbox*}
